%% kapa: exam=1 round=37 year=2026 period=1 form=A subject=real_estate subject_ko=부동산학원론 seq=7 number=87
%% source: Q-Net 감정평가사 기출문제 (cst00309) artlSeq=5258924 file=raw/1008/5258924/2026년 제37회 감정평가사 1차 1교시.pdf
%% conversion: kordoc 4.24.51 -> tools/gen1.py (auto); review=proofread-sonnet; flags=none
\begin{question}{87}
부동산등기제도에 관한 설명으로 옳은 것은?
\begin{choices}
  \item 임차권과 권리질권은 등기할 수 없다.
  \item 「부동산등기법」상 예고등기제는 현재 시행되고 있다.
  \item 부동산의 등기기록에는 표제부, 갑구, 을구, 병구를 둔다.
  \item 등기권리의 순위는 등기기록 중 다른 구(區)에서 한 등기 상호간에는 순위번호에 따른다.
  \item 구분건물(區分建物)의 등기기록은 표제부가 2개로서 1동 건물의 표제부와 전유부분의 표제부가 있다.
\end{choices}
\end{question}
\answer{5}
